\section{Policy Instruments and Economic Feedback}
\label{sec:policy-feedback}

\subsection{Carbon price}

A carbon price changes the relative cost of technologies in CLEWS, favouring those that
emit less. In the economy, a carbon tax is also a payment to government, and the use of
the revenue matters: returning it to households has different effects from spending it or
reducing debt. OG--CLEWS keeps the resource cost of cleaner technologies separate from the
payment of the tax, so that each is counted once and matched to where it goes.

In the Philippine scenario, a tax of US\$50 per tonne of CO$_2$ on
household energy, with the revenue returned to households as transfers, lowers GDP by about
0.05 percent over the first decade and by 0.02 percent in the long run. Consumption falls at
first and then recovers (Figure~\ref{fig:alone}).

\subsection{Investment incentives}

Policy can also change the incentive to invest. A tax credit, accelerated depreciation or a
lower business tax reduces the cost of investing in generation. It does not make the physical
asset cheaper; it shifts part of its cost from the investor to the government. Firms respond
by investing more, and the public budget bears the cost. This is the instrument that draws
capital into the sector, in contrast with the change in technology described above. The two
can be used together when the change in technology and the policy support are specified
separately.

\subsection{Planning discount rate}

CLEWS compares costs that occur at different dates using a discount rate. A lower rate gives
more weight to future savings, and favours assets with high initial costs and low operating
costs, such as renewable generation. OG-Core produces the return that the economy itself
requires on capital, which reflects how households choose between consuming today and in
the future. In the Philippine example, this return averages about 8.2 percent over the first
decade. At that rate, a cost of 100 twenty years from now is worth about 21 today; at 4
percent it is worth about 46. The choice of rate therefore changes which technologies look
cheapest. The economy's return can inform that choice, but it is not automatically the right
rate for public planning: risk, taxes and the social perspective must also be taken into
account.

\subsection{Demand}

A resource plan is designed to meet a given demand, but the economy's response changes that
demand. Higher electricity costs reduce consumption and production, while investment and
better health can increase activity. OG--CLEWS returns the change in demand to CLEWS, which
can then revise the capacity and investment the plan requires. In the Philippine
scenario, electricity production is about 6.6 percent below the baseline in 2030, when the price increase peaks, and about 3.5 percent below it in 2040. A
plan that ignores this response risks building more capacity than the economy will use.
